% TOP_PROBLEM: {"id":30007753,"problem_number":"LOCAL-30007753","title":"Coordinating multiple tax instruments","category_id":21,"category":{"id":21,"name":"economics","display_name":"Economics","description":"Open economic theory statements, published research agendas, and proposed scoped specifications; question provenance and historical priority are qualified per record.","slug":"economics","order_index":21,"created_at":"2026-10-08T00:00:00Z"},"status":"research_question","external_url":null,"legacy_ids":[],"published":false,"statement":"In the explicitly specified setting: Which combinations of income, consumption, capital, and family taxes can improve welfare without shifting distortions elsewhere? The mathematical statement above fixes the target, assumptions and scope of this catalogue formulation.","economics":{"original_id":242,"field":"Taxation and social insurance","kind":"design","formalization_basis":"adapted_research_question","source_status":"research_agenda","priority_status":"earliest_found","priority_note":"Earliest explicit agenda verified in this audit, not a claim of first historical publication. The mathematical scope is authored for this catalogue.","earliest_verified_reference":{"authors":["Louis Kaplow"],"title":"Optimal Income Taxation","year":2022,"url":"https://laweconcenter.law.harvard.edu/wp-content/uploads/2024/11/Kaplow_1083.pdf","doi":"10.3386/w30199","locator":"Olin Discussion Paper 1083, July 2022; sections 3.2, 7, 8.2 and 9, printed pp. 89–90","evidence_summary":"The conclusion identifies further work on heterogeneous preferences, family composition, endogenous wages, and integrating tax instruments. This is an explicit research agenda, not the exact optimization model below."},"current_reference":{"authors":["Louis Kaplow"],"title":"Optimal Income Taxation","year":2024,"url":"https://www.aeaweb.org/articles?id=10.1257/jel.20221647","doi":"10.1257/jel.20221647","locator":"Journal of Economic Literature 62(2), 637–738; publisher abstract","evidence_summary":"The published survey describes heterogeneous preferences, families, endogenous wages and interactions with other instruments."},"formalization_note":"Catalogue-authored scoped formalization. Population, instruments, welfare objective and identification restrictions must be instantiated before claiming a resolution; source authors are not credited with these added assumptions."},"statusQualification":"Published question or research agenda verified at the cited locator. The mathematical specification is adapted for this catalogue and includes additional assumptions; source publication does not establish first-ever priority or certify this exact model remains unsolved. Catalogue-authored scoped formalization. Population, instruments, welfare objective and identification restrictions must be instantiated before claiming a resolution; source authors are not credited with these added assumptions.","statusReviewedAt":"2026-10-08"}
% FIRST_POSTED: 2026-10-08
% ENTRY_KIND: statement_only
% !TeX program = lualatex
\documentclass[11pt]{article}
\usepackage{fontspec}
\usepackage[a4paper,margin=25mm]{geometry}
\usepackage{amsmath,amssymb,mathtools}
\usepackage{xurl}
\usepackage[hidelinks]{hyperref}
\newcommand{\sref}[2]{\hyperlink{source:#1:#2}{[S#2]}}
\setlength{\emergencystretch}{3em}
\begin{document}
% problemId:30007753
\section*{Coordinating multiple tax instruments}\label{30007753}
\subsection{Definitions and mathematical statement}
Consider a specified economy of households with private types \(\theta\sim F\), firms and a government. A policy \(p=(T_y,T_k,t_c,T_f)\) consists of labor-income and capital-income schedules, a vector of commodity tax rates, and transfers conditioned on observable family characteristics. Specify bounded admissible schedules \(\mathcal P\), preferences, technologies, information restrictions and administrative cost \(C(p)\). For policy \(p\), let \(\mathcal E(p)\) be the nonempty set of competitive equilibria including endogenous earnings, savings, prices and family choices. For equilibrium \(e\), define social welfare \(W(e)=\int G_\theta(U_\theta(e))\,dF(\theta)\) using explicitly fixed utility indices and welfare functions. Define net government receipts \(B(p,e)\) as all tax receipts less transfers and \(C(p)\). The formal task is to characterize
\[p^*\in\arg\max_{p\in\mathcal P}\inf_{e\in\mathcal E(p)} W(e)\quad\text{subject to}\quad B(p,e)\geq R\ \text{for all }e\in\mathcal E(p),\]
where \(R\) is a specified funding requirement. Establish implementable sufficient conditions for a revenue-neutral joint reform to improve welfare over a specified baseline, accounting for cross-instrument behavioral and price effects. Distinguish changes in actual allocations from relabeling the same effective tax wedge. Identify restrictions under which a smaller instrument set attains the same value and quantify any loss when those restrictions fail.

\par\textbf{Formulation and scope.} Catalogue-authored scoped formalization. Population, instruments, welfare objective and identification restrictions must be instantiated before claiming a resolution; source authors are not credited with these added assumptions.

\par\textbf{Open-question provenance.} An explicit question or research agenda was verified at the source locator. This does not by itself certify that every later version remains unresolved \sref{30007753}{1}.

\par\textbf{Historical priority.} Earliest explicit agenda verified in this audit, not a claim of first historical publication. The mathematical scope is authored for this catalogue.

\subsection{Short English statement}
In the explicitly specified setting: Which combinations of income, consumption, capital, and family taxes can improve welfare without shifting distortions elsewhere? The mathematical statement above fixes the target, assumptions and scope of this catalogue formulation.

\subsection{Sources}
\begin{itemize}\raggedright
\item[\textnormal{[S1]}]\hypertarget{source:30007753:1}{} Louis Kaplow. 2022. Optimal Income Taxation. Olin Discussion Paper 1083, July 2022; sections 3.2, 7, 8.2 and 9, printed pp. 89–90.
\url{https://laweconcenter.law.harvard.edu/wp-content/uploads/2024/11/Kaplow_1083.pdf}.
DOI: 10.3386/w30199.
{\small\textit{Earliest verified question/agenda reference; historical priority qualified below. The conclusion identifies further work on heterogeneous preferences, family composition, endogenous wages, and integrating tax instruments. This is an explicit research agenda, not the exact optimization model below.}}

\item[\textnormal{[S2]}]\hypertarget{source:30007753:2}{} Louis Kaplow. 2024. Optimal Income Taxation. Journal of Economic Literature 62(2), 637–738; publisher abstract.
\url{https://www.aeaweb.org/articles?id=10.1257/jel.20221647}.
DOI: 10.1257/jel.20221647.
{\small\textit{Supporting or current literature; see evidence qualification. The published survey describes heterogeneous preferences, families, endogenous wages and interactions with other instruments.}}
\end{itemize}
\end{document}
